\scanpage{82}
% manuscript line 1
$Z$ constructible dans $X$ quasi compact, \inserted{dont} tt pt a un syst. fond. de voisinages rétrocompacts dans $X$.
% manuscript line 2
\deleted{$Z$ loc constructible}
% manuscript line 3
\[
\left\{\begin{aligned}
&Z\text{ quasi cpct et}\\
&\forall x\in Z,\ \exists\text{ un voisinage }U_x\text{ de }x\text{ dans }X,\text{ tel que }Z\cap U_x\text{ constr. dans }U_x\\
&\text{(alors tt ouvert }U'\subset U,\ U'\ni x,\text{ a la }=\text{ propriété)}.
\end{aligned}\right.
\]
% manuscript line 6
Nécessité évidente [car $Z$ rétrocompact dans $X$, \deleted{\unsure{donc}} quasi cpct et on peut prendre $U=X$

% manuscript line 8
Suffisance. $Z$ est recouvert par un nb fini des $U_x$, \unsure{soient|soit} $U_i$. les $U_i\cap Z$ sont constructibles dans $U_i$, donc dans $X$, donc $Z=\bigcup Z_i$ est constr.

% manuscript line 11
\uline{Corollaire} Supposons $Z$ loc. fermé. Alors $Z$ \deleted{loc.} constructible ssi $Z$ est \uline{quasi compact}, et tt pt $x\in Z$ a un voisinage ouvert $U$ dans $X$ tel que $Z\cap U$ \deleted{soit fermé} \deleted{\unsure{comm|et}} \uline{constructible} dans $U$, i.e. $U-Z\cap U$ rétro-\deleted{\unsure{cp.}}ct dans $U$ [i.e. $U-Z\cap U$ quasi-compact, si on \inserted{prend $U$ quasi-cpct, et si on} suppose que dans $X$, l'intersection de deux ouverts quasi cpcts est quasi cpcte).

% manuscript line 18
\uline{Exemple} Si $X$ est un préschéma quasi-séparé quasi-compact (p. ex. un schéma quasi-compact, tel un schéma affine) \deleted{\unsure{Alors}} $Z$ loc. fermé dans $X$, alors $Z$ constructible $\Longleftrightarrow$ $Z$ quasi cpct, \uuline{et} tt $x\in Z$ a un voisinage ouvert affine $U$ tel que $Z\cap U$ soit défini par un \uline{nb fini} d'\uline{équations}
% archive sheet number: 81

% Reverse-face typescript shows through here. See scan 83.
